\documentclass[11pt,a4paper]{article}
\usepackage[margin=24mm]{geometry}
\usepackage{amsmath,amssymb,booktabs,microtype}
\usepackage[hidelinks]{hyperref}
\newcommand{\dd}{\mathrm d}
\newcommand{\cA}{\mathcal A}
\newcommand{\cC}{\mathcal C}
\title{Route G2-R: Joint Output Modes and Detector Recoil}
\author{Conditional tree-level predictions from the unchanged G2 action}
\date{23 September 2026}
\begin{document}
\maketitle
\begin{abstract}
We turn the existing G2 collision amplitudes into joint mode--recoil
predictions, keeping all open channels and distinguishing scattering
events from mode-conversion events. A mass-only detector cannot resolve
the sign of its compact momentum. For nonelastic channels in the
specified sharp-beam preparation, outgoing ordinary momentum can
disambiguate that sign. This is a model-dependent reconstruction, not
an independent proof of compact-momentum conservation. No new action,
parameter fit, absolute encounter probability or external drive is added.
\end{abstract}

\section{Inputs and two different event ensembles}
The G2 interaction remains the positive local potential
$\kappa_5|\Phi|^2|X|^2$ on $\mathbb M^{3,1}\times S^1$.
Its normalized circle overlap gives
\begin{equation}
 m_n^2=M_5^2+(n+\alpha)^2/R^2,\quad
 \mu_l^2=M_D^2+l^2/R^2,\quad
 \mathcal M=-g\,\delta_{n+l,m+j},\quad g=\kappa_5/(2\pi R).
\end{equation}
Keep the entire G2 engineering card:
$R=1$, $M_5=0.5$, $\alpha=0.25$, $M_D=5$, $g=0.05$,
incoming $\Phi$ label $n=1$, and ordinary center-of-mass momentum
$p_i=0.2$. The detector coefficients are exactly finite,
\begin{equation}
 c_l=\mathcal N e^{-l^2/(4w^2)}e^{-\mathrm i l\theta_0},
 \quad l=-4,\ldots,4,\quad w=1.2,\quad\theta_0=0.7,\quad
 w_l=|c_l|^2,\quad\sum_l w_l=1.
\end{equation}
These are prepared-state inputs, not empirical fitted values.
For each incoming $l$, set
\begin{align}
 \sqrt{s_l}&=\sqrt{m_n^2+p_i^2}+\sqrt{\mu_l^2+p_i^2},\quad
 j=n+l-m,\\
 p_{lm}&=\frac{\sqrt{[s_l-(m_m+\mu_j)^2]
                             [s_l-(m_m-\mu_j)^2]}}{2\sqrt{s_l}},\\
 K_{lm}&=\sigma_{lm}v_{{\rm M},l}
 =\frac{g^2p_{lm}}{16\pi E_nE_l\sqrt{s_l}}.
 \label{eq:rate}
\end{align}
Here $K_{lm}=0$ unless $\sqrt{s_l}>m_m+\mu_j$; closed channels
must be rejected before evaluating the radical. The distinct-particle
phase-space convention is the same as G2 and the PDG review
\cite{pdg}. Define
\begin{equation}
 \cA=\{(l,m):w_l>0,\ \sqrt{s_l}>m_m+\mu_{n+l-m}\},\qquad
 \cC=\{(l,m)\in\cA:m\ne n\}.
\end{equation}
$\cA$ contains all \emph{scattering events}, including elastic
$m=n$ collisions. $\cC$ contains only mode-changing events.
Neither includes the unscattered identity amplitude as a collision.
This distinction is required before assigning a normalized fraction.

\section{Joint distribution, not a single target fraction}
For $B=\cA$ or $\cC$, let
\begin{equation}
 Z_B=\sum_{(l,m)\in B}w_lK_{lm},\qquad
 \boxed{P_B(m,j)=
 \frac{w_{m+j-n}K_{m+j-n,m}}{Z_B}
 \mathbf 1_{\{(m+j-n,m)\in B\}}}.
 \label{eq:joint}
\end{equation}
For $Z_B>0$ these dimensionless quantities sum to one and are nonnegative.
They are probabilities \emph{conditional on the selected type of
scattering event}, in the common-overlap dilute approximation of G2.
They are not the probability that an arbitrary incident packet scatters.
For unequal encounter overlaps $I_l$, replace $w_lK_{lm}$ by
$w_lI_lK_{lm}$ and renormalize. No values of $I_l$ are invented here.
The common coupling $g^2$ cancels from \eqref{eq:joint}, but not
from the event rate.

The marginal and conditional predictions are
\begin{equation}
 P_B^\Phi(m)=\sum_jP_B(m,j),\quad
 P_B^X(j)=\sum_mP_B(m,j),\quad
 P_B(j\mid m)=\frac{P_B(m,j)}{P_B^\Phi(m)}
\end{equation}
where the denominator is nonzero. The complete table, rather than
only $m=0$, is saved in the numerical report. Conditional on each
incoming $l$, support lies on $m+j=n+l$. A packet with several $l$
therefore need not exhibit a single perfectly anticorrelated line.

Event conditioning changes the incoming detector distribution to
\begin{equation}
 P_B^{\rm in}(l)=\frac{w_l\sum_{m:(l,m)\in B}K_{lm}}{Z_B}.
\end{equation}
Conservation identities must use this distribution:
\begin{align}
 \langle m\rangle_B+\langle j\rangle_B-\langle l\rangle_B&=n,\\
 \operatorname{Var}_B(l)
 &=\operatorname{Var}_B(m)+\operatorname{Var}_B(j)
        +2\operatorname{Cov}_B(m,j),\\
 \langle j-l\mid m\rangle_B&=n-m .
\end{align}
Comparing selected outgoing recoil with the \emph{unconditioned}
incident mean would mix momentum transfer with event-selection bias.
Perfect anticorrelation applies to the transfers
$\Delta j=j-l=-\Delta m=n-m$, so
$\operatorname{Cov}(\Delta m,\Delta j)=-\operatorname{Var}(\Delta m)$.
It does not generally apply to the output labels $m,j$ themselves.

For the declared card there are 25 open $(l,m)$ channels.
The total and mode-changing rate coefficients are respectively
$Z_{\cA}=1.85661445\,10^{-6}$ and $Z_{\cC}=1.63742848\,10^{-6}$.
The all-scattering output marginal is
\begin{center}
\begin{tabular}{rrrrr}
\toprule
Output $m$ & $-2$ & $-1$ & $0$ & $1$ (elastic)\\
\midrule
$P_{\cA}^{\Phi}(m)$ & $0.00785985$ & $0.27119405$ & $0.60288929$ & $0.11805680$\\
\bottomrule
\end{tabular}
\end{center}
Thus $Z_{\cC}/Z_{\cA}=0.88194320$ of scattering events change
mode under this prescription. This is not an incident conversion efficiency.

For fixed $(m,j)$, $l=m+j-n$ is fixed, and the ideal narrow-beam
momentum-angle law is
\begin{equation}\label{eq:lines}
 \dd P_B(m,j,p,\Omega)
 =P_B(m,j)\,\delta(p-p_{lm})\,\dd p\,\frac{\dd\Omega}{4\pi}.
\end{equation}
The two ordinary outgoing momenta are $p_{lm}\widehat{\boldsymbol u}$
and $-p_{lm}\widehat{\boldsymbol u}$ in the common COM frame.
Isotropy is a tree contact-interaction prediction. Finite beam
widths and instrumental response would broaden these ideal lines.

\section{What a mass-only detector actually resolves}
The detector mass $\mu_j=\sqrt{M_D^2+j^2/R^2}$ measures
$r=|j|$, not its sign. The corresponding joint distribution is
\begin{equation}
 \overline P_B(m,r)=
 \begin{cases}
 P_B(m,0),&r=0,\\
 P_B(m,r)+P_B(m,-r),&r>0 .
 \end{cases}
\end{equation}
This is an incoherent sum over orthogonal recoil states. Degeneracy
of their measured masses does not make their amplitudes add.
Where both signs have support,
\begin{equation}
 P_B(j\mid m,r)=\frac{P_B(m,j)}{\overline P_B(m,r)},\qquad |j|=r .
\end{equation}
For $\alpha=0.25$ the $\Phi$ masses determine their signed mode
labels uniquely in the ideal tree spectrum: an equality for distinct
$m,m'$ would require $m+m'=-2\alpha$, impossible for integers.
At $2\alpha\in\mathbb Z$ this step can fail because of $\Phi$
degeneracies. G1 already requires eigenspace, not arbitrary basis-state,
interpretations at such degeneracies.

Only diagonal recoil event rates are calculated. They are independent
of $\theta_0$, and an incoming detector mixture diagonal in $l$ with
the same $w_l$ gives the same rates. Thus this joint table is not an
entanglement witness or evidence that the preparation was localized.
It also does not construct a real measurement apparatus for signed $j$.

\section{A kinematic way to resolve the sign in conversion events}
Suppose the ideal readout provides $(m,r,p_f)$, the parameters and
incoming $p_i$ are known, and the event is in the COM frame.
The measured final energy and inferred incoming detector momentum are
\begin{align}
 E_f&=\sqrt{m_m^2+p_f^2}
        +\sqrt{M_D^2+r^2/R^2+p_f^2},\\
 E_n&=\sqrt{m_n^2+p_i^2},\\
 l^2&=R^2\left[(E_f-E_n)^2-M_D^2-p_i^2\right].
\end{align}
Write $d=m-n$ so that $l=d+j$. If $d\ne0$,
\begin{equation}\label{eq:reconstruct}
 l^2=d^2+2dj+r^2
 \quad\Longrightarrow\quad
 \boxed{j=\frac{l^2-d^2-r^2}{2d}} .
\end{equation}
The result must be consistent with $j=\pm r$, integer mode labels,
the prepared support and the open-channel condition. This is
\emph{model-dependent kinematic reconstruction}: the dispersion law
and momentum conservation were used to obtain it. Checking that the
reconstructed $j$ obeys the same conservation law is not an independent
test of that law. The independently comparable predictions are the
joint recoil line positions, event fractions and angular distribution.

The uniqueness can also be proved without inversion. For the two
candidate signs $j=\pm r$, the initial labels are $l=d\pm r$, and
\begin{equation}
 \sqrt{s_+}-\sqrt{s_-}
 =\frac{4dr/R^2}
 {\sqrt{M_D^2+p_i^2+(d+r)^2/R^2}
  +\sqrt{M_D^2+p_i^2+(d-r)^2/R^2}}.
\end{equation}
This is nonzero whenever $dr\ne0$. At fixed final masses,
\begin{equation}
 \frac{\dd E_f}{\dd p_f}
 =p_f\left(\frac1{E'_\Phi}+\frac1{E'_X}\right)>0
\end{equation}
above threshold, so both open signs have different ordinary
outgoing momenta. For elastic $d=0$, $l=j$ and $p_f=p_i$;
these measurements cannot recover the sign. For $r=0$, the
only possible value is already $j=0$.

\subsection*{A concrete two-line prediction}
For the selected output $m=0$, detector mass $\mu_1=5.09901951$
corresponds to $r=1$. Two allowed components give
\begin{center}
\begin{tabular}{rrrrr}
\toprule
Incoming $l$ & Outgoing $j$ & $\sqrt{s_l}$ & $p_f$ & $P(j\mid m=0,r=1)$\\
\midrule
$0$ & $+1$ & $6.36506416$ & $1.02175458$ & $0.77457541$\\
$-2$ & $-1$ & $6.74994319$ & $1.36191317$ & $0.22542459$\\
\bottomrule
\end{tabular}
\end{center}
These conditional weights use the same prepared packet and common
overlap, with equal ideal acceptance. The mass-only readout merges
the lines; the ideal momentum readout
separates them by $\Delta p_f=0.34015860$ in the common energy unit.
For a deterministic bound $|p_{\rm meas}-p_{\rm true}|\leq\epsilon_p$,
with all other inputs treated as exact, a sufficient disjoint-interval
criterion is $2\epsilon_p<\Delta p_f$. This is a required resolution,
not an assumed experimental error bar. Uncertain incoming momentum,
parameters, mass assignment or frame must also be propagated before
using that criterion experimentally.

\section{Readout boundaries and a bounded next step}
A calibrated classical readout response
$\mathcal R(b\mid m,j,p,\Omega)$ would give
\begin{equation}
 P_B^{\rm obs}(b)=\sum_{m,j}P_B(m,j)
 \int\frac{\dd\Omega}{4\pi}\,
 \mathcal R(b\mid m,j,p_{m+j-n,m},\Omega).
\end{equation}
Include a missed-event outcome to make $\sum_b\mathcal R=1$.
Conditioning on detected events requires renormalization and may
change fractions. Only the ideal signed-label and mass-only readouts
are evaluated here; no detector efficiency, angular cut or noise
distribution is fitted or invented.

The executable imports the unchanged G2 card, enumerates every open
channel, and verifies normalization, marginals, Bayes identities,
event-conditioned recoil, angular normalization, gauge relabeling,
sign folding and reconstruction. Reports preserve the code/input
hashes and the complete joint matrices. The underlying G2 tree-EFT
and common-overlap limitations remain. In particular no interacting
loop error or absolute four-dimensional packet encounter probability
is supplied.

The next useful choice is a specified finite-resolution readout to
test whether these lines survive realistic input spread and response.
If an absolute encounter probability is needed instead, specify
normalized four-dimensional packets and use the existing G2 scattering
amplitude. An energy-accounted driven background remains an unused
backup. Neither option demands a new mass fit, lattice scan or action.

\begin{thebibliography}{9}
\bibitem{pdg}
Particle Data Group, \emph{Kinematics}, Review of Particle Physics
(2024), sections 49.5--49.5.1,
\href{https://pdg.lbl.gov/2024/reviews/rpp2024-rev-kinematics.pdf}{official review}.
\end{thebibliography}
\end{document}
